\documentclass[11pt]{article}
\usepackage{amsmath,amssymb,amsthm}
\usepackage{algorithm}
\usepackage{algpseudocode}
\usepackage{geometry}
\usepackage{xcolor}
\geometry{margin=1in}
\newtheorem{theorem}{Theorem}
\newtheorem{lemma}{Lemma}
\newtheorem{remark}{Remark}
\newcommand{\grad}{\nabla}
\newcommand{\R}{\mathbb{R}}

\begin{document}

\section*{Heavy‑Ball Momentum (HB) Method}

The method under study updates the iterate $x_k\in\R^d$ by adding a momentum term that depends on the \emph{previous iterate} (not on a look‑ahead gradient).  
The update reads
\[
\boxed{\;x_{k+1}=x_k+\beta\,(x_k-x_{k-1})-\alpha\,\grad f(x_k)\;},
\qquad k\ge 1,
\]
where $\alpha>0$ is a step‑size (learning rate) and $\beta\in[0,1)$ is the momentum parameter.  
The algorithm is often called the \textbf{Heavy‑Ball} method.

---------------------------------------------------------------------

\section*{Mathematical Formulation}
\begin{align}
\text{Given: }&\;x_0\in\R^d,\;x_{-1}=x_0,\;\alpha>0,\;\beta\ge 0. \nonumber\\[2mm]
\text{For }k=0,1,2,\dots:&\nonumber\\
\qquad\text{(1) Compute gradient: } & g_k:=\grad f(x_k). \nonumber\\
\qquad\text{(2) Momentum term: } & m_k:=\beta\,(x_k-x_{k-1}). \nonumber\\
\qquad\text{(3) Update: } & x_{k+1}=x_k+m_k-\alpha g_k . \nonumber
\end{align}

---------------------------------------------------------------------

\section*{Pseudocode}
\begin{algorithm}[H]
\caption{Heavy‑Ball Momentum}
\begin{algorithmic}[1]
\Require Objective $f:\R^d\!\to\!\R$, step size $\alpha>0$, momentum $\beta\ge0$, initial point $x_0$.
\State $x_{-1}\gets x_0$
\For{$k=0,1,2,\dots$}
    \State $g_k\gets \grad f(x_k)$
    \State $m_k\gets \beta\,(x_k-x_{k-1})$
    \State $x_{k+1}\gets x_k+m_k-\alpha\,g_k$
\EndFor
\State \textbf{output} $x_{K}$ (or the best iterate seen)
\end{algorithmic}
\end{algorithm}

---------------------------------------------------------------------

\section*{Dry Run Example}
Consider the one‑dimensional quadratic
\[
f(w)=(w-3)^2 .
\]
Then $\grad f(w)=2(w-3)$.  
Choose $\alpha=0.1$, $\beta=0.5$, and initialise $w_0=0$, $w_{-1}=0$.

\begin{center}
\begin{tabular}{c|c|c|c}
\textbf{Iter. $k$} & $g_k=2(w_k-3)$ & $m_k=\beta (w_k-w_{k-1})$ & $w_{k+1}=w_k+m_k-\alpha g_k$ \\ \hline
0 & $-6$ & $0$ & $w_1=0+0-0.1(-6)=0.6$ \\
1 & $2(0.6-3)=-4.8$ & $0.5(0.6-0)=0.3$ & $w_2=0.6+0.3-0.1(-4.8)=1.38$ \\
2 & $2(1.38-3)=-3.24$ & $0.5(1.38-0.6)=0.39$ & $w_3=1.38+0.39-0.1(-3.24)=2.104$ 
\end{tabular}
\end{center}

Objective values:
\[
f(w_0)=9,\; f(w_1)=(0.6-3)^2=5.76,\; f(w_2)=(1.38-3)^2\approx2.62,\; f(w_3)=(2.104-3)^2\approx0.80 .
\]
The iterates clearly move toward the minimiser $w^\star=3$.

---------------------------------------------------------------------

\newpage
\section*{Assumptions}
\begin{enumerate}
\item \textbf{Strong convexity.} $f$ is $\mu$‑strongly convex ($\mu>0$):
\begin{equation}
f(y)\ge f(x)+\langle\grad f(x),y-x\rangle+\frac{\mu}{2}\|y-x\|^2,\qquad\forall x,y\in\R^d .
\label{eq:strong}
\end{equation}
Equivalently,
\begin{equation}
\langle\grad f(x)-\grad f(y),x-y\rangle\ge \mu\|x-y\|^2 .
\label{eq:strong-grad}
\end{equation}

\item \textbf{Smoothness.} $f$ has $L$‑Lipschitz gradient ($L\ge\mu$):
\begin{equation}
\|\grad f(x)-\grad f(y)\|\le L\|x-y\|,\qquad\forall x,y\in\R^d .
\label{eq:smooth}
\end{equation}
Equivalently,
\begin{equation}
\langle\grad f(x)-\grad f(y),x-y\rangle\le L\|x-y\|^2 .
\label{eq:smooth-grad}
\end{equation}

\item \textbf{Parameter region.} The step size $\alpha$ and momentum $\beta$ satisfy
\begin{equation}
0<\alpha<\frac{2}{L+\mu},\qquad
0\le\beta\le\bigl(1-\sqrt{\alpha\mu}\,\bigr)^2 .
\label{eq:params}
\end{equation}
Under \eqref{eq:params} the spectral radius of the linearised error recursion is strictly smaller than one (see Theorem below).

\item \textbf{Existence of a unique minimiser.} Because $f$ is strongly convex, there exists a unique $x^\star\in\R^d$ such that $\grad f(x^\star)=0$.
\end{enumerate}

---------------------------------------------------------------------

\section*{Main Convergence Theorem}
\begin{theorem}[Linear convergence of Heavy‑Ball]\label{thm:linear}
Assume \ref{eq:strong}–\ref{eq:params}.  Let $\{x_k\}$ be generated by the Heavy‑Ball update.  Define the error $e_k:=x_k-x^\star$.  Then for all $k\ge 0$
\begin{equation}
\|e_k\|\;\le\;\rho^{\,k}\,\|e_0\| ,
\qquad\text{where }\;
\rho:=\max\Bigl\{\;|1-\alpha\mu+\beta|,\;|1-\alpha L+\beta|\;\Bigr\}<1 .
\label{eq:linear-rate}
\end{equation}
Consequently,
\[
f(x_k)-f(x^\star)\;\le\;\frac{L}{2}\,\rho^{\,2k}\,\|e_0\|^2 .
\]
\end{theorem}

---------------------------------------------------------------------

\section*{Theoretical Properties}
\begin{itemize}
\item \textbf{Stability.} The admissible region \eqref{eq:params} guarantees that the linear operator governing the error dynamics is a contraction; hence the method is stable for any starting point.

\item \textbf{Hyper‑parameter sensitivity.} The bound $\beta\le(1-\sqrt{\alpha\mu})^2$ shows that a larger step size $\alpha$ forces a smaller momentum.  Setting $\beta=0$ reduces the method to plain gradient descent, whose rate is $|1-\alpha\mu|$.  For the optimal choice $\alpha=\frac{4}{(\sqrt{L}+\sqrt{\mu})^2}$ and $\beta=\bigl(\frac{\sqrt{L}-\sqrt{\mu}}{\sqrt{L}+\sqrt{\mu}}\bigr)^2$, one obtains the classical accelerated rate $\rho=\frac{\sqrt{L}-\sqrt{\mu}}{\sqrt{L}+\sqrt{\mu}}$.

\item \textbf{Comparison to baselines.} Compared with Gradient Descent (GD) the Heavy‑Ball method attains a faster linear factor $\rho_{\text{HB}}<\rho_{\text{GD}}$ whenever $\beta>0$.  However, unlike Nesterov’s Accelerated Gradient, Heavy‑Ball does not enjoy a universal optimal worst‑case bound; the analysis here is tight for quadratics.

\end{itemize}

---------------------------------------------------------------------

\section*{Discussion}
The theorem shows that under strong convexity and smoothness the Heavy‑Ball method converges \emph{geometrically} to the unique minimiser.  The proof proceeds by bounding the error recursion with a two‑dimensional linear system whose spectral radius can be expressed explicitly in terms of $\alpha,\beta,\mu,L$.  The parameter region \eqref{eq:params} is exactly the set where this spectral radius is $<1$.  The result matches the classic analysis of Polyak (1964) and demonstrates that the algorithm is robust to the choice of initial point.

---------------------------------------------------------------------

\newpage
\section*{Preliminaries (Definitions and Lemmas)}
\begin{lemma}[Gradient Lipschitz equivalence]\label{lem:grad-bounds}
For a $\mu$‑strongly convex and $L$‑smooth function $f$, the Jacobian matrix $H(x)$ defined by the mean‑value theorem satisfies
\[
\mu I \;\preceq\; H(x) \;\preceq\; L I ,\qquad\forall x\in\R^d .
\]
Consequently,
\[
\mu\|v\|^2 \le v^\top H(x) v \le L\|v\|^2 ,\qquad\forall v\in\R^d .
\]
\end{lemma}
\begin{proof}
By the integral form of the mean‑value theorem,
\[
\grad f(x)-\grad f(y)=\Bigl(\int_{0}^{1} \nabla^2 f\bigl(y+t(x-y)\bigr)\,dt\Bigr)(x-y)
=:H_{xy}(x-y).
\]
Strong convexity \eqref{eq:strong-grad} yields $v^\top H_{xy}v\ge\mu\|v\|^2$ and smoothness \eqref{eq:smooth-grad} yields $v^\top H_{xy}v\le L\|v\|^2$.  Since $H_{xy}$ is a convex combination of Hessians, the same bounds hold for any $H(x)$ defined at a single point.  ∎
\end{proof}

---------------------------------------------------------------------

\section*{Main Proof (Step‑by‑Step Derivation)}
We prove Theorem~\ref{thm:linear}.  Throughout the proof we write $e_k:=x_k-x^\star$ and use $g_k:=\grad f(x_k)$.  Recall that $\grad f(x^\star)=0$.

\subsection*{Step 1.  Error recursion}
\begin{align}
e_{k+1}
&=x_{k+1}-x^\star \nonumber\\
&=x_k+\beta(x_k-x_{k-1})-\alpha g_k - x^\star \nonumber\\
&=e_k+\beta(e_k-e_{k-1})-\alpha g_k . \label{eq:err-rec}
\end{align}
[Step 1] is a direct substitution of the update rule and the definition of $e_k$.

\subsection*{Step 2.  Linearisation of the gradient}
By Lemma~\ref{lem:grad-bounds} there exists a symmetric matrix $H_k$ with
\[
\mu I\preceq H_k\preceq L I
\]
such that
\[
g_k=\grad f(x_k)=\grad f(x^\star)+H_k e_k = H_k e_k .
\tag{2}
\]
[Step 2] follows from the mean‑value theorem applied between $x_k$ and $x^\star$.

\subsection*{Step 3.  Substitute (2) into (1)}
\begin{equation}
e_{k+1}= \bigl(1+\beta-\alpha H_k\bigr)e_k -\beta e_{k-1}. \label{eq:lin-rec}
\end{equation}
[Step 3] uses $g_k=H_k e_k$.

\subsection*{Step 4.  2‑dimensional state representation}
Define the stacked vector
\[
z_k:=\begin{bmatrix}e_k\\ e_{k-1}\end{bmatrix}\in\R^{2d}.
\]
Then \eqref{eq:lin-rec} can be written as a linear system
\[
z_{k+1}= \underbrace{\begin{bmatrix}
(1+\beta)I-\alpha H_k & -\beta I\\[1mm]
I & 0
\end{bmatrix}}_{=:\,\mathcal{M}_k}\,z_k .
\tag{3}
\]
[Step 4] is straightforward matrix algebra.

\subsection*{Step 5.  Bounding the spectral radius}
Because $H_k$ is symmetric and satisfies $\mu I\preceq H_k\preceq L I$, the eigenvalues of $\mathcal{M}_k$ lie in the interval between the eigenvalues of the two extreme matrices
\[
\mathcal{M}_\mu:=\begin{bmatrix}
(1+\beta-\alpha\mu)I & -\beta I\\[1mm]
I & 0
\end{bmatrix},
\qquad
\mathcal{M}_L:=\begin{bmatrix}
(1+\beta-\alpha L)I & -\beta I\\[1mm]
I & 0
\end{bmatrix}.
\tag{4}
\]
Indeed, for any eigenpair $(\lambda,v)$ of $\mathcal{M}_k$ we have
\[
\lambda v = \mathcal{M}_k v
\;\Longrightarrow\;
\lambda^2 - (1+\beta-\alpha \theta)\lambda +\beta =0,
\]
where $\theta$ is an eigenvalue of $H_k$ and therefore $\theta\in[\mu,L]$.  Hence the characteristic polynomial of $\mathcal{M}_k$ is
\[
p_{\theta}(\lambda)=\lambda^2-(1+\beta-\alpha\theta)\lambda+\beta .
\tag{5}
\]
[Step 5] uses the block‑structure of $\mathcal{M}_k$ and the fact that $I$ commutes with $H_k$.

\subsection*{Step 6.  Roots of the polynomial}
The two roots of \eqref{eq:5} are
\[
\lambda_{\pm}(\theta)=\frac{1+\beta-\alpha\theta\pm\sqrt{(1+\beta-\alpha\theta)^2-4\beta}}{2}.
\tag{6}
\]
Their magnitudes are maximised at the endpoints $\theta=\mu$ or $\theta=L$ because the discriminant is a concave quadratic in $\theta$.  Therefore the spectral radius of $\mathcal{M}_k$ is bounded by
\[
\rho:=\max\bigl\{\,|\lambda_{\pm}(\mu)|,\;|\lambda_{\pm}(L)|\,\bigr\}.
\tag{7}
\]

\subsection*{Step 7.  Explicit expression of $\rho$}
Observe that for any $\theta\in\{\mu,L\}$ the polynomial $p_\theta$ satisfies
\[
p_\theta(1)=1-(1+\beta-\alpha\theta)+\beta = \alpha\theta- \beta .
\]
Since $0\le\beta\le(1-\sqrt{\alpha\mu})^2$ (assumption \eqref{eq:params}), we have $p_\theta(1)\ge0$ and $p_\theta(0)=\beta>0$.  Moreover,
\[
p_\theta\bigl(1-\alpha\theta\bigr)= (1-\alpha\theta)^2-(1+\beta-\alpha\theta)(1-\alpha\theta)+\beta
= \beta(1-\alpha\theta)\ge0 .
\]
Thus the two real roots lie in the interval $[\,1-\alpha\theta,\;1-\alpha\theta+\beta\,]$.  Consequently,
\[
|\lambda_{\pm}(\theta)|\le \max\bigl\{|1-\alpha\theta|,\;|1-\alpha\theta+\beta|\bigr\}.
\]
Taking the worst case over $\theta\in\{\mu,L\}$ yields exactly the constant $\rho$ defined in \eqref{eq:linear-rate}:
\[
\rho=\max\Bigl\{|1-\alpha\mu+\beta|,\;|1-\alpha L+\beta|\Bigr\}<1 .
\tag{8}
\]
The strict inequality follows from the parameter restriction \eqref{eq:params} (see e.g. Polyak 1964).

\subsection*{Step 8.  Contraction of the state}
From \eqref{eq:3} we have $\|z_{k+1}\|\le \|\mathcal{M}_k\|\,\|z_k\|$.  Using the induced 2‑norm and the bound $\|\mathcal{M}_k\|\le\rho$ that follows from (8), we obtain by induction
\[
\|z_k\|\le\rho^{\,k}\|z_0\|.
\tag{9}
\]
Since $\|z_k\|^2=\|e_k\|^2+\|e_{k-1}\|^2\ge\|e_k\|^2$, (9) implies
\[
\|e_k\|\le\rho^{\,k}\|e_0\|,
\]
which is exactly the error bound claimed in \eqref{eq:linear-rate}.  This completes the proof of the first part of Theorem~\ref{thm:linear}.  ∎

\subsection*{Step 9.  Function‑value rate}
Using $L$‑smoothness,
\[
f(x_k)-f(x^\star)\le \frac{L}{2}\|e_k\|^2
\le \frac{L}{2}\rho^{\,2k}\|e_0\|^2 .
\]
Thus the function values converge with rate $\rho^{2k}$, as stated.  ∎

---------------------------------------------------------------------

\section*{Rate Derivation (With Constants)}
For the optimal choice
\[
\alpha^\star=\frac{4}{(\sqrt{L}+\sqrt{\mu})^{2}},\qquad
\beta^\star=\Bigl(\frac{\sqrt{L}-\sqrt{\mu}}{\sqrt{L}+\sqrt{\mu}}\Bigr)^{2},
\]
one computes
\[
\rho^\star = \frac{\sqrt{L}-\sqrt{\mu}}{\sqrt{L}+\sqrt{\mu}} \;<\;1 .
\]
Hence
\[
\|x_k-x^\star\|\le \Bigl(\tfrac{\sqrt{L}-\sqrt{\mu}}{\sqrt{L}+\sqrt{\mu}}\Bigr)^{k}\|x_0-x^\star\|,
\qquad
f(x_k)-f(x^\star)\le \frac{L}{2}\Bigl(\tfrac{\sqrt{L}-\sqrt{\mu}}{\sqrt{L}+\sqrt{\mu}}\Bigr)^{2k}\|x_0-x^\star\|^{2}.
\]

---------------------------------------------------------------------

\section*{Special Cases}
\begin{enumerate}
\item \textbf{No momentum ($\beta=0$).}  The update reduces to gradient descent and $\rho=|1-\alpha\mu|$, which recovers the classic GD linear rate.

\item \textbf{Quadratic functions.}  For $f(x)=\tfrac12 x^\top Q x-b^\top x$ with $Q$ having eigenvalues in $[\mu,L]$, the matrix $H_k$ equals $Q$ for all $k$.  The analysis becomes exact, and the bound $\rho$ is the true spectral radius of the error dynamics.

\item \textbf{Limit $\alpha\to 0$.}  When the step size tends to zero, $\rho\to 1-\beta$; thus a large momentum $\beta\approx 1$ yields very slow contraction, reflecting the need for a sufficiently large step size to make progress.
\end{enumerate}

---------------------------------------------------------------------

\end{document}